\section{A faithful labelled heat source and the original theta endpoints}
\label{sec:labelled-heat-original-theta}

\paragraph{LH1. The received obstruction specifies a construction.}
The supplied continuation \emph{Untitled 1793} proposes the shifted
Gaussian sum with its zero mode retained. We retain that proposal and
prove its domains, its source maps and an exact return to the original
zeta function. No assertion that every heat operation is already an
endomorphism of the arithmetic semiring is used.
Let \(L\) be any nontrivial bounded distributive lattice, with no
finiteness assumption, and retain
\[
 S=G_L(\mathbb Z)
 =\{(n,1_L):n\in\mathbb Z\}\cup\{(0,\lambda):\lambda\in L\},
 \quad \tau=(0,0_L),\quad e=(0,1_L).             \tag{LH1}
\]
There is one copy of \((0,1_L)\). The coefficient operations are
\((a,\lambda)+(b,\mu)=(a+b,\lambda\vee\mu)\) and
\((a,\lambda)(b,\mu)=(ab,\lambda\wedge\mu)\). Carrier membership
holds because a nontop label forces zero amplitude; distributivity
follows from ring distributivity and the distributive lattice law.
The supported-zero ideal is prime: it is precisely
\(\{(0,\lambda):\lambda\in L\}\), since multiplying \(e\)
by these same elements gives all of them. A product belongs to
it exactly when its integer amplitudes have zero product, hence
one amplitude is zero. It is proper because it excludes \((1,1_L)\).
This proves the required prime directly without deleting any label.

The phrase that both zero elements must map to zero under every
multiplicative observation is too strong. The amplitude map
\(p:S\to\mathbb R\), \(p(n,1)=n\), \(p(0,\lambda)=0\),
is a unital semiring map and does identify the whole zero fibre.
But every unital multiplicative map to \(\mathbb R\) sends
\(e\) to an idempotent, hence either zero or one. If it sends
\(e\) to one, the equations \(e(n,1)=e\) force the image of
every supported integer to be one. This case occurs: the support
character of \(G_{\{0,1\}}(\mathbb Z)\) sends all supported
integers to one and \(\tau\) to zero, multiplicatively. It is
not additive into the real ring, since \(1+(-1)=e\) would give
\(2=1\). For every additive map to a characteristic-zero vector
space, instead, \(z_\lambda+z_\lambda=z_\lambda\) forces the
image of \(z_\lambda\) to zero. These are exact restrictions on
the respective maps, not an irrelevance argument about the zero fibre.

\paragraph{LH2. A source which keeps every labelled point.}
Let \(V_L=\bigoplus_{\lambda\in L}\mathbb C v_\lambda\), with
independent basis vectors, and let \(\mathcal D'_L\) be the direct
sum of tempered-distribution spaces indexed by \(L\). Finite
linear combinations of the following embedded points form a
subspace \(\mathcal A_L\):
\[
 \iota(n,1_L)=\delta_n v_{1_L},\qquad
 \iota(0,\lambda)=\delta_0v_\lambda\quad(\lambda\ne1_L).
                                                               \tag{LH2}
\]
This is an injective map of sets, and its linear extension is an
isomorphism from the free complex vector space on the set \(S\)
to \(\mathcal A_L\). To prove independence, separate the lattice
coordinates. In the top coordinate the finitely many Dirac masses
are independent: smooth compactly supported functions in disjoint
small intervals around their integers extract each coefficient.
The other coordinates each contain one independent Dirac mass.
This free construction retains the original coefficient operations
as the bilinear point maps
\[
 \begin{split}
 m_+(\delta_a v_\lambda\otimes\delta_bv_\mu)
    &=\delta_{a+b}v_{\lambda\vee\mu},\\
 m_\times(\delta_a v_\lambda\otimes\delta_bv_\mu)
    &=\delta_{ab}v_{\lambda\wedge\mu}.
 \end{split}\tag{LH3}
\]
They are defined on the displayed source bases, extended linearly
on \(\mathcal A_L\otimes\mathcal A_L\). On the original embedded
points they reproduce every original operation, including the
additive identity \(\tau\), the multiplicative unit and both
zero elements. We do not assert that ordinary vector addition
reproduces the original semiring addition; the proved map is LH3.
For example \(\iota(e)\) is a nonzero vector although \(e+e=e\).

\paragraph{LH3. Positive-time heat with an exact inverse on its image.}
For actual heat time \(T>0\) define
\[
 k_T(x)=(4\pi T)^{-1/2}\exp[-x^2/(4T)],\qquad
 \mathcal H_T U=k_T*U                                      \tag{LH4}
\]
coordinatewise. Differentiation gives
\(\partial_Tk_T=\partial_x^2k_T\), including the derivative
of the complete prefactor. The Gaussian Fourier transform, in the
original convention \(\widehat f(\xi)=\int f(x)e^{-2\pi i x\xi}dx\),
is
\[
 \widehat k_T(\xi)=e^{-4\pi^2T\xi^2}.              \tag{LH5}
\]
For a direct proof, integration by parts yields
\(I'(\xi)=-8\pi^2T\xi I(\xi)\); the Gaussian integral gives
\(I(0)=1\). Solving this first-order equation proves LH5.
Completing the square in the absolutely convergent convolution
integral, retaining both prefactors, gives
\(k_T*k_U=k_{T+U}\).

Convolution by this Schwartz function is defined on each tempered
distribution and has Fourier multiplier LH5. It is injective.
Indeed if \(e^{-4\pi^2T\xi^2}\widehat U=0\), test locally
against a compactly supported smooth function \(\phi\), after
dividing that test by the smooth nonzero multiplier. This gives
\(\langle\widehat U,\phi\rangle=0\). Compact-test uniqueness
implies \(\widehat U=0\) and hence \(U=0\). Thus the inverse
on the actual range is the distributional operation
\[
 U=\mathcal F^{-1}
       \left(e^{4\pi^2T\xi^2}\widehat{\mathcal H_TU}\right).
                                                               \tag{LH6}
\]
This inverse is asserted on this range; multiplying an arbitrary
tempered distribution by the growing exponential need not yield
a tempered distribution. No unrestricted backward heat assertion
is made. Every finite support label is retained separately.

In particular \(\mathcal H_T\iota(e)=k_Tv_{1_L}\) and
\(\mathcal H_T\iota(\tau)=k_Tv_{0_L}\) are different nonzero
vectors for every positive time. The exact original arithmetic
operations on their heat images are
\[
 m_{+,T}=\mathcal H_Tm_+
        (\mathcal H_T^{-1}\otimes\mathcal H_T^{-1}),\qquad
 m_{\times,T}=\mathcal H_Tm_\times
        (\mathcal H_T^{-1}\otimes\mathcal H_T^{-1}).             \tag{LH7}
\]
Their domain is \(\mathcal H_T\mathcal A_L\otimes
\mathcal H_T\mathcal A_L\), and their output is
\(\mathcal H_T\mathcal A_L\). Substitution gives, for example,
\(m_{+,T}(k_T(x-a)v_\lambda,k_T(x-b)v_\mu)
=k_T(x-a-b)v_{\lambda\vee\mu}\), and the product receiver
has \(ab\) and \(\lambda\wedge\mu\). This supplies explicit
faithful heat images of the semiring points and their operations.
It does not identify these operations with pointwise addition
and multiplication of heat functions.

\paragraph{LH4. The actual integer comb and its marked zero.}
The infinite integer comb \(D=\sum_{n\in\mathbb Z}\delta_n\)
is tempered, because every Schwartz function has summable values
on the integers. Split its source as the retained pair
\((D-\delta_0,\delta_0)\), whose sum is \(D\). With a finite
set of lower labels \(J\subseteq L\setminus\{1_L\}\), retain
the source
\[
 \mathbf D_J=Dv_{1_L}+\delta_0\sum_{\lambda\in J}v_\lambda.
                                                               \tag{LH8}
\]
No infinite sum of unrelated lattice coordinates is needed. The
directed union over finite \(J\) retains every individual lower
label. Heat gives exactly
\[
 \mathcal H_T\mathbf D_J(x)
  =\sum_{n\in\mathbb Z}k_T(x-n)v_{1_L}
        +k_T(x)\sum_{\lambda\in J}v_\lambda.       \tag{LH9}
\]
The separate marked top zero term is \(k_T(x)v_{1_L}\).
It satisfies the same heat equation and has mass one for all
\(T>0\). Thus this construction moves the marked zero profile
without erasing it. Applying the inverse LH6 componentwise to the
retained pair recovers that pair and every lower coordinate. The
sum alone does not specify an arbitrary choice of marked summand.

The top coordinate in LH9 is periodic, but the separate marked
Gaussian is not periodic. Periodizing that separate Gaussian
gives the top circle heat function and thereby forgets which
integer was originally marked. Keeping the pair avoids this
additional loss. This identifies a precise limitation of scalar
circle periodization while retaining its connecting map.

\paragraph{LH5. Shifted theta and its two arithmetic sheets.}
Set \(x=(4\pi T)^{-1}\) and use a separate displacement
variable \(b\in\mathbb R\). Then the exact top heat kernel is
\[
 \sum_{n\in\mathbb Z}k_T(b-n)
   =\sqrt{x}\,K(x;b),\qquad
 K(x;b)=\sum_{n\in\mathbb Z}e^{-\pi(n+b)^2x}.      \tag{LH10}
\]
Changing \(n\) to \(-n\) proves the sign in this equality.
Direct differentiation or LH4 proves the heat equation in
\((T,b)\); \(b\) is displacement, not the Rodgers--Tao time.
On compact \(x>0\) ranges every derivative converges by a
Gaussian majorant. The original marked top zero and each lower
coordinate have the term \(E(x;b)=e^{-\pi b^2x}\). The exact
labelled comparison is
\[
 \big((K-E)v_{1_L},\ E v_{1_L},\ (E v_\lambda)_{\lambda\in J}\big),
                                                               \tag{LH11}
\]
with the specified sum map to LH9 after multiplication by
\(\sqrt x\). All components remain available before that sum.

For \(0<b<1\), \(\Re s>1\), termwise Mellin integration yields
\[
 \begin{split}
 \int_0^\infty K(x;b)x^{s/2-1}dx
  =\pi^{-s/2}\Gamma(s/2)
    \left[b^{-s}+\sum_{n\ge1}(n+b)^{-s}
                    +\sum_{n\ge1}(n-b)^{-s}\right].
 \end{split}\tag{LH12}
\]
To prove interchange, replace \(s\) by its real part in the
absolute value, integrate each positive Gaussian by
\(y=\pi(n+b)^2x\), and use convergence of the resulting
inverse-power series. The two last sums are exactly
\(\zeta(s,1+b)\) and \(\zeta(s,1-b)\) on this half-plane.
Their named functions are not substitutes for either full series
or the displayed marked term \(b^{-s}\).

At \(b=0\) the marked term is exactly one as a function of
\(x\), and its integral over \((0,\infty)\) has no common
convergence half-plane. The correct retained pair has the
unmarked Mellin identity
\[
 \int_0^\infty(K(x;0)-1)x^{s/2-1}dx
      =2\pi^{-s/2}\Gamma(s/2)\zeta(s),\quad\Re s>1,
                                                               \tag{LH13}
\]
alongside the separate marked constant and all lower labels.
Consequently it recovers the original \(\zeta\) by multiplying
by \(\pi^{s/2}/[2\Gamma(s/2)]\). This is an equality with
every factor retained, not an identification of the two functions.

\paragraph{LH6. Both endpoint terms, derived from the actual Gaussian.}
For completeness the Gaussian Poisson identity needed here follows
without a formal deletion of the zero term. The smooth periodic
sum \(K(x;b)\) has Fourier coefficient at integer \(m\) equal to
\(x^{-1/2}e^{-\pi m^2/x}\): unfold the sum over a unit interval
and apply the Gaussian Fourier calculation LH5. Its absolutely
convergent Fourier series therefore gives
\[
 K(x;b)=x^{-1/2}\sum_{m\in\mathbb Z}
                 e^{-\pi m^2/x}e^{2\pi i m b}.     \tag{LH14}
\]
At zero displacement set \(\psi(x)=\sum_{n\ge1}e^{-\pi n^2x}\).
The zero Fourier coefficient and the zero source coefficient give
separately
\(\psi(x)=x^{-1/2}\psi(1/x)+(x^{-1/2}-1)/2\).
Split the Mellin integral of \(\psi\) at one and substitute
\(y=1/x\) only on \((0,1)\). Direct integration of each of the
two endpoint terms then proves
\[
 \zeta(s)=\frac{\pi^{s/2}}{\Gamma(s/2)}
 \left[\frac1{s-1}-\frac1s+
 \int_1^\infty\psi(x)
       \left(x^{s/2-1}+x^{-(s+1)/2}\right)dx\right].             \tag{LH15}
\]
The integral is entire in \(s\), since the Gaussian decay on
\([1,\infty)\) controls all powers of \(\log x\) on compact
sets. Thus LH15 extends meromorphically and retains both endpoint
terms. It agrees with the original-source theta formula used by
\href{https://arxiv.org/abs/1801.05914v5}{Rodgers and Tao}, while
CF1--29 supply the exact subsequent differential and line-module
map. The supplied continuation's proposed shifted sum is thereby
proved as an exact scalar Mellin representation, with the faithful
labelled distribution source retained by LH11. The scalar shifted
sum itself has \(K(x;b)=K(x;1-b)\) for \(0<b<1\), by the
integer change of index, so it does not determine the displacement
on that whole interval. Keeping \(E(x;b)=e^{-\pi b^2x}\) resolves
this ambiguity there: for any fixed \(x>0\) this is strictly
decreasing in \(b>0\), with inverse
\(b=\sqrt{-\log E(x;b)/(\pi x)}\).

\paragraph{LH7. The obstruction and its computed consequence.}
The actual lower labels do not enter the scalar series in LH15:
projection onto the top coordinate has kernel precisely the direct
sum of the lower-coordinate distribution spaces. Its restriction
to the linear span of \(Dv_{1_L}\) and
\(\{\delta_0v_\lambda:\lambda\in J\}\) has the finite kernel spanned by
\(\delta_0v_\lambda\), \(\lambda\in J\), whose heat images are
the nonzero \(k_Tv_\lambda\). Keeping LH9 and LH11 makes the
source and its heat evolution faithful; projecting loses exactly
that identified subspace. The labelled heat generator acts there
by \(\partial_x^2\), which is nonzero on its Gaussian images.
Thus this is an explicit alternative positive-time heat construction
with an active marked zero mode. No assertion that its positivity
or its ordinary circle spectrum settles the original RH is part
of the calculation. The maps LH2--15 specify which original data
it retains and how its scalar part returns to the original zeta.
